% Compile with LuaLaTeX
\documentclass[11pt,a4paper]{article}


% ==========================================
% 1. MODERN FONTS & LANGUAGES
% ==========================================
\usepackage{fontspec}
% Tip: For Romance languages and IPA, "Charis SIL" or "Doulos SIL" are the best fonts. 
% If you have them installed, uncomment the next line:
% \setmainfont{Charis SIL} 

\usepackage{polyglossia}
\setmainlanguage{english} % change to german if needed
% Load the Romance languages you study
\setotherlanguages{french,spanish,italian,portuguese,romanian,latin}

% ==========================================
% 1b. TYPOGRAPHY & QUOTES
% ==========================================
\usepackage{microtype}   % Better line breaks, especially helpful for multilingual docs
\usepackage{csquotes}    % Language-aware quotes: \enquote{...} picks the right marks
\usepackage{booktabs}
\usepackage{graphicx}

% ==========================================
% 2. LINGUISTICS & PHONETICS
% ==========================================
\usepackage{tipauni}     % Modern, LuaLaTeX-safe IPA symbols
\usepackage{gb4e}        % For interlinear glosses (standard in linguistics)
\noautomath              % CRITICAL: Prevents gb4e from crashing with modern packages
\usepackage{forest}      % For drawing modern syntax trees

% ==========================================
% 3. MATH & LAYOUT
% ==========================================
\usepackage{amsmath,amssymb,mathtools}
\usepackage{geometry}
\geometry{margin=2.5cm}
\usepackage{hyperref}
\usepackage{xcolor}
\usepackage{siunitx}

% ==========================================
% 4. HYPERSETUP
% ==========================================
 
\hypersetup{
    colorlinks=true,
    linkcolor=blue,
    urlcolor=blue
}

% ==========================================
% 5.    CODE LISTINGS
% ==========================================

\usepackage{listings}
\lstset{
    basicstyle=\ttfamily\small,
    breaklines=true,
    frame=single,
    backgroundcolor=\color{gray!10}
}

\title{A Modern Guide to \textbf{LuaLaTeX} \\ \large For Linguists and Beginners (2026 Edition)}
\author{Adapted for Modern TeX Live}
\date{\today}

\begin{document}

\maketitle
\tableofcontents
\vspace{1em}

\section{Introduction: What is LuaLaTeX?}

In the past, most people used \textbf{pdfLaTeX}. However, pdfLaTeX was designed in the 1990s, before the modern internet and Unicode. Today, if you want to use the actual fonts installed on your computer (like Times New Roman, Arial, or specialized linguistic fonts) and type any character from any language directly (like \textbf{ə, ʃ, ɛ̃}), you need a modern engine.

\textbf{LuaLaTeX} is that modern engine. It has three major advantages:
\begin{enumerate}
    \item \textbf{Native Unicode Support:} You can type characters from any language (French, German, IPA) directly into your document. You don't need weird commands like \texttt{\textbackslash'e} for é; you just type é.
    \item \textbf{Modern Fonts:} It uses the \texttt{fontspec} package to access TrueType and OpenType fonts installed on your Mac or PC.
    \item \textbf{Lua Scripting:} It embeds the Lua programming language, allowing for advanced automation (though you rarely need this as a beginner).
\end{enumerate}

\section{The Golden Rules: Switching from pdfLaTeX to LuaLaTeX}

If you find a tutorial online that was written for pdfLaTeX, \textbf{do not copy its preamble blindly}. There are four golden rules for LuaLaTeX:

\begin{itemize}
    \item \textbf{Never use \texttt{inputenc}.} LuaLaTeX handles UTF-8 natively. (Delete \texttt{\textbackslash usepackage\{inputenc\}}).
    \item \textbf{Do not load \texttt{fontenc} or \texttt{textcomp}.} Under LuaLaTeX, \texttt{fontenc} is superseded by \texttt{fontspec}, and \texttt{textcomp} is already folded into current LaTeX. If you see them in an older tutorial, just delete those lines.
    \item \textbf{Use \texttt{fontspec} instead.} This is how you load fonts.
    \item \textbf{Use \texttt{polyglossia} instead of \texttt{babel}.} \texttt{polyglossia} is the modern language package built for Unicode.
\end{itemize}

\section{Migration Table: Pre-2015 vs 2026}

If you are porting a template from an older tutorial, this is the shortest path. Everything on the left is from pre-2015 tutorials you may still find on StackExchange, Overleaf, or university handouts; everything on the right is the current practice under LuaLaTeX.

\begin{center}
\small
\renewcommand{\arraystretch}{1.6}  % add some air between the lines
\begin{tabular}{p{1.4cm} p{4.6cm} p{4.3cm} p{3cm}}
\toprule
\textbf{Topic} & \textbf{Old way (pre-2015)} & \textbf{New way (2026)} & \textbf{Remark} \\
\midrule
Input encoding & \texttt{\textbackslash usepackage[utf8]\{inputenc\}} & Delete it. & LuaLaTeX is UTF-8 native. \\
Font encoding & \texttt{\textbackslash usepackage[T1]\{fontenc\}} & \texttt{\textbackslash usepackage\{fontspec\}} & \texttt{textcomp} is also obsolete. \\
Fonts & \texttt{\textbackslash usepackage\{times\}} and friends & \texttt{\textbackslash setmainfont\{...\}} & Any OpenType font on your OS. \\
Languages & \texttt{\textbackslash usepackage[french]\{babel\}} & \texttt{\textbackslash usepackage\{polyglossia\}} & Then \texttt{\textbackslash setotherlanguages\{...\}}. \\
IPA & \texttt{\textbackslash usepackage\{tipa\}} & \texttt{\textbackslash usepackage\{tipauni\}}\,$^{\ast}$ & See note below. \\
Quotes & \texttt{``...''}, \texttt{<<...>>} by hand & \texttt{\textbackslash enquote\{...\}} (\texttt{csquotes}) & Language-aware. \\
Glosses & \texttt{\textbackslash usepackage\{gb4e\}} & \texttt{\textbackslash usepackage\{gb4e\}} + \texttt{\textbackslash noautomath} & Same package, essential workaround. \\
Compile & \texttt{pdflatex}, click twice & \texttt{latexmk -lualatex -pdf} & Handles multi-pass + Biber. \\
\bottomrule
\end{tabular}
\end{center}

\noindent $^{\ast}$ \texttt{tipauni} is a Unicode-era replacement for \texttt{tipa}. Because it uses true Unicode characters, it is compatible with \texttt{fontspec} and \texttt{polyglossia}; the old \texttt{tipa} is not.

\section{How to Handle Fonts}

Because LuaLaTeX uses your operating system's fonts, you use \texttt{\textbackslash setmainfont} to choose your text font. By default, LuaLaTeX uses \textit{Latin Modern}, which is fine.

If you want to change it (for example, to a font that has excellent IPA support, like \textit{Charis SIL} or \textit{Doulos SIL}), you simply add this to your preamble:
\begin{lstlisting}
\usepackage{fontspec}
\setmainfont{Charis SIL} % The font must be installed on your computer!
\end{lstlisting}

\section{Modern Linguistics Packages}

In the 2010s, linguists used the \texttt{tipa} package for phonetic symbols. However, \texttt{tipa} breaks easily with modern fonts. Today, we use \textbf{\texttt{tipauni}} (TIPA Unicode), which is specifically designed for LuaLaTeX.

\subsection{IPA Symbols}
Load the package like this:
\begin{lstlisting}
\usepackage{tipauni}
\end{lstlisting}
Then, you can use the familiar \texttt{\textbackslash textipa\{\}} shortcuts, or simply type the Unicode symbols directly if your font supports it.

\subsection{Glosses and Examples (\texttt{gb4e})}
The \texttt{gb4e} package is still the standard for numbered examples and interlinear glosses. However, it sometimes conflicts with math and phonetic packages.
\textbf{The Fix:} Always put \texttt{\textbackslash noautomath} right after loading \texttt{gb4e} to prevent crashes.

\begin{lstlisting}
\usepackage{gb4e}
\noautomath % Fixes a known infinite loop bug
\end{lstlisting}

\section{A Complete Minimal Preamble}

Here is a safe, modern preamble you can copy for your computational linguistics assignments:

\begin{lstlisting}
\documentclass[11pt,a4paper]{article}

% 1. Modern Engine & Fonts
\usepackage{fontspec}
% \setmainfont{Charis SIL} % Uncomment to change font

% 2. Languages
\usepackage{polyglossia}
\setmainlanguage{english}
\setotherlanguages{french,spanish,german}

% 3. Math
\usepackage{amsmath,amssymb}

% 4. Linguistics
\usepackage{tipauni}     % Modern IPA
\usepackage{gb4e}
\noautomath              % CRITICAL for gb4e to work with other packages
\usepackage{forest}      % Syntax trees

% 5. Layout
\usepackage{geometry}
\geometry{margin=2.5cm}
\usepackage{hyperref}
\end{lstlisting}

\section{Common Troubleshooting (The "Gotchas")}

\subsection{Common First Errors — a quick lookup}

Scan this table before anything else. Nine out of ten first-time errors are in it.

\begin{center}
\small
\renewcommand{\arraystretch}{1.8}  % add some air between the lines
\begin{tabular}{p{4.8cm} p{3.6cm} p{5.5cm}}
\toprule
\textbf{Symptom / Error message} & \textbf{Likely cause} & \textbf{Fix} \\
\midrule
\texttt{Unicode character ... not set up for use with LaTeX} & The current font lacks that glyph. & Switch to a font with wide coverage: \texttt{\textbackslash setmainfont\{Charis SIL\}} or Doulos SIL. \\
Missing references, broken Table of Contents & First pass has not yet resolved labels. & Compile again, or use \texttt{latexmk -lualatex -pdf}. \\
\texttt{File 'fontspec.sty' not found} & You are on a minimal TeX installation. & Install the recommended TeX Live collections (see installation section). \\
Strange IPA output or missing symbols & Legacy \texttt{tipa} workflow combined with a non-tipa font. & Switch to \texttt{tipauni} and a Unicode-capable font. \\
\texttt{gb4e} crashes on a subscript or superscript & Missing \texttt{\textbackslash noautomath}. & Add \texttt{\textbackslash noautomath} immediately after \texttt{\textbackslash usepackage\{gb4e\}}. \\
\texttt{luaotfload: font ... not found} & Font cache is stale. & Linux: \texttt{fc-cache -fv}. macOS: \texttt{luaotfload-tool --update}. \\
\bottomrule
\end{tabular}
\end{center}

\subsection{Missing Packages (\texttt{tlmgr})}
If you get an error saying \texttt{File 'xyz.sty' not found}, it means your TeX installation is missing a package. If you are using a minimal installation (like BasicTeX on Mac), you must install it via the terminal:
\begin{lstlisting}
sudo tlmgr install ifmtarg
\end{lstlisting}

\subsection{Stale Font Caches}
If you install a new font (such as \textit{Charis SIL}) and LuaLaTeX cannot find it, the font database needs to be refreshed.

On \textbf{Ubuntu / Linux}, after dropping the \texttt{.ttf} files into \texttt{\textasciitilde/.local/share/fonts/}:
\begin{lstlisting}
fc-cache -fv
\end{lstlisting}

On \textbf{macOS} the equivalent one-off command is:
\begin{lstlisting}
luaotfload-tool --update
\end{lstlisting}
Do \textbf{not} prepend \texttt{sudo} to \texttt{luaotfload-tool} --- it writes to your user cache and must run as your normal user. (On Windows the system handles font installation automatically; a LuaLaTeX cache refresh is rarely needed.)

\subsection{Compilation Requirements}
Unlike pdfLaTeX, LuaLaTeX sometimes requires \textbf{two compilation passes} to get the Table of Contents and cross-references right. If you are using \texttt{biblatex}, the full sequence is: LuaLaTeX \(\rightarrow\) Biber \(\rightarrow\) LuaLaTeX \(\rightarrow\) LuaLaTeX.

Clicking the compile button four times is tedious and error-prone. Instead, open a terminal in the project folder and run:
\begin{lstlisting}
latexmk -lualatex -pdf your-document.tex
\end{lstlisting}
\texttt{latexmk} inspects the output after each pass, detects whether references are still unresolved, runs Biber when the \texttt{.bib} file has changed, and recompiles exactly as many times as needed. It ships with every modern TeX Live installation; there is nothing to install.

For reproducibility (so a document still compiles cleanly a few years from now), the usual recipe is \texttt{latexmk} + \texttt{git} + a pinned TeX Live version. For teaching material or theses that may outlive their author's laptop, a Docker image such as \texttt{texlive/texlive} wraps the whole toolchain and can rebuild the exact same PDF on any machine.

\section{A Note on Lua (Why We Use Python Instead)}

LuaLaTeX embeds the Lua programming language. You can write loops, read CSV files, and do calculations directly inside your \texttt{.tex} source using \texttt{\textbackslash directlua}. However, as a computational linguist, your real data processing almost certainly belongs in Python or R --- not inside the typesetting engine. We include one tiny example here so you understand the bridge exists, but we deliberately do not teach Lua programming in this kit.

\begin{lstlisting}
\documentclass{article}
\usepackage{fontspec}

\begin{document}

\directlua{
  local words = {"chat", "chien", "maison"}
  tex.sprint("Corpus contains " .. #words .. " items.")
}

\end{document}
\end{lstlisting}

When compiled, this prints ``Corpus contains 3 items.'' in the body of the PDF. The Lua block runs at compile time, builds a string, and \texttt{tex.sprint} injects it into the LaTeX output as if you had typed it yourself.

For anything substantial, do your corpus processing in Python or R, write the result to a \texttt{.csv} or a \texttt{.tex} fragment, and let LaTeX \texttt{\textbackslash input} that file. That is a cleaner division of labour than embedding complex logic inside your document --- and it keeps your analysis reproducible with the tools your discipline already uses.

\section{Recommended Learning Resources}

\begin{itemize}
    \item \textbf{Overleaf's 30-Minute Intro:} \url{https://www.overleaf.com/learn/latex/Learn_LaTeX_in_30_minutes} (Remember to change the compiler to LuaLaTeX in the Overleaf menu!)
    \item \textbf{LearnLaTeX.org:} \url{https://www.learnlatex.org/} (Highly recommended interactive tutorial).
    \item \textbf{Detexify:} \url{http://detexify.kirelabs.org/classify.html} (Draw a symbol with your mouse, and it tells you the LaTeX command).
\end{itemize}

\end{document}